\title{DeepSeekMath: Pushing the Limits of Mathematical Reasoning in Open Language Models}

\begin{abstract}
The task of mathematical reasoning presents considerable difficulties for language models, stemming from its inherent complexity and structured requirements. This study presents DeepSeekMath~7B, a model developed by extending the pre-training of DeepSeek-Coder-Base-v1.5 7B using a substantial dataset containing 120B tokens related to mathematics, curated from Common Crawl alongside data encompassing both natural language and code. Remarkably, DeepSeekMath~7B attains a score of 51.7\% on the highly challenging, competition-grade MATH benchmark without leveraging any supplementary toolkits or ensemble methods, narrowing the gap with models such as Gemini-Ultra and GPT-4. When evaluated via self-consistency using 64 samples, DeepSeekMath~7B reaches a performance of 60.9\% on MATH. The advancements in mathematical reasoning evident in DeepSeekMath~are attributed to two central innovations: (1) an advanced data filtering and selection pipeline that extracts maximum value from freely available web resources, and (2) the implementation of Group Relative Policy Optimization (GRPO)—an adaptation of the Proximal Policy Optimization (PPO) algorithm—which simultaneously improves mathematical reasoning and makes PPO more memory-efficient.
\end{abstract}

\section{Introduction}

Large language models (LLM) have dramatically transformed how artificial intelligence tackles mathematical reasoning, leading to significant improvements on quantitative reasoning benchmarks \citep{MATH} and geometry reasoning benchmarks \citep{trinh2024solving}. These models are also proving valuable as tools that aid individuals in working through difficult mathematical tasks \citep{tao}. Nevertheless, premier systems such as GPT-4 \citep{gpt4} and Gemini-Ultra \citep{gemini} remain closed-source, while the open-source alternatives currently available exhibit a substantial performance deficit in comparison.

This paper presents DeepSeekMath, a specialized language model designed for mathematical reasoning, which demonstrates superior mathematical performance relative to open-source alternatives and comes close to matching GPT-4 on established academic evaluations. The foundation of this advancement is the DeepSeekMath Corpus—a large, high-quality pre-training dataset comprising 120B math-related tokens. This corpus is assembled from Common Crawl (CC) using a fastText-based classifier \citep{joulin2016fasttext}. For the initial training phase, the classifier utilizes samples from OpenWebMath \citep{openwebmath} as positive cases, alongside a wide array of general web pages serving as negatives. The trained classifier is subsequently applied to the CC to extract additional likely positive samples, which are further validated and refined through manual annotation. Incorporating this curated data, the classifier undergoes additional training to enhance its accuracy. Evaluations show the efficacy of this corpus, as evidenced by our DeepSeekMath-Base 7B model’s results: achieving 64.2\% accuracy on GSM8K \citep{gsm8k} and 36.2\% on the MATH competition dataset \citep{MATH}, surpassing the performance of Minerva 540B \citep{minerva}. The DeepSeekMath Corpus is also multilingual, leading to observed improvements on Chinese mathematical benchmark tasks \citep{wei2023cmath,agieval}. We suggest that our methodology for mathematical data curation offers a valuable foundation for further work in the field, with considerable opportunity for future advancements.

DeepSeekMath-Base is built upon DeepSeek-Coder-Base-v1.5 7B \citep{deepseek-coder}, given our observation that initializing from a code-focused pretrained model yields better outcomes than commencing with a generic LLM. In addition, we find that mathematical pre-training not only strengthens mathematical proficiency but also bolsters overall reasoning, as reflected by improvements on the MMLU \citep{mmlu} and BBH \citep{bbh} benchmarks.

Post pre-training, we further fine-tune DeepSeekMath-Base with mathematical instruction data that incorporates chain-of-thought reasoning \citep{cot}, program-of-thought methodologies \citep{pot,pal}, and tool-enhanced reasoning tasks \citep{tora}. The final model, DeepSeekMath-Instruct 7B, surpasses all other 7B-scale models and achieves results on par with open-source instruction-tuned models at the 70B parameter level.

Moreover, we propose Group Relative Policy Optimization (GRPO), a novel reinforcement learning (RL) approach based on Proximal Policy Optimization (PPO) \citep{schulman2017proximal}. GRPO eliminates the need for a critic model and estimates baselines using scores aggregated within groups, thus greatly reducing the computational requirements for training. By leveraging only a segment of English instruction tuning datasets, GRPO drives marked enhancements over DeepSeekMath-Instruct, both for in-domain problems (GSM8K: 82.9\% $\rightarrow$ 88.2\%, MATH: 46.8\% $\rightarrow$ 51.7\%) and for out-of-domain tasks such as CMATH (84.6\% $\rightarrow$ 88.8\%) during RL optimization. We also present a unified conceptual framework connecting methods including Rejection Sampling Fine-Tuning (RFT) \citep{yuan2023scaling}, Direct Preference Optimization (DPO) \citep{dpo}, PPO, and GRPO, revealing that these can all be regarded as variants or simplifications of RL strategies. Through comprehensive experiments, such as contrasting online and offline training, evaluating outcome- versus process-oriented supervision, and comparing single-turn to iterative RL, we systematically probe the core aspects of this framework. Finally, we articulate the mechanisms by which our RL approach enhances the capabilities of instruction-tuned models and outline promising avenues for more advanced RL, grounded in our unified perspective.

\subsection{Contributions}
This work advances scalable approaches to math pre-training and systematically investigates the role of reinforcement learning.

\textbf{Math Pre-Training at Scale}

\begin{itemize}[topsep=0pt]
    \item We demonstrate that the openly available Common Crawl corpus holds significant utility for mathematical data extraction. Through a carefully crafted selection pipeline, we assemble the DeepSeekMath Corpus, which comprises 120B tokens derived from web sources curated specifically for mathematical relevance. This dataset substantially surpasses previous collections, being nearly 7 times larger than that of Minerva \citep{minerva} and exceeding the scale of OpenWebMath \citep{openwebmath} by a factor of 9.

    \item Our foundational model, DeepSeekMath-Base 7B, achieves performance on par with the Minerva 540B model \citep{minerva}. These results suggest that raw parameter count alone does not dictate a model’s ability to perform mathematical reasoning; instead, models of smaller scale but trained on more targeted and higher-quality data can deliver competitive results.

    \item Experimental results regarding the training regimen are also shared. We find that initializing models with code-focused pre-training prior to math-focused training enhances their competence in mathematical problem-solving—both in scenarios with tool assistance and those without. This observation lends evidence to the debated idea that code pre-training boosts reasoning capacity, particularly in mathematical contexts.

    \item Despite widespread practice of incorporating arXiv papers for math model training, we observe that doing so yields negligible benefit across the range of mathematical benchmarks employed in this study.
\end{itemize}

\textbf{Exploration and Analysis of Reinforcement Learning}

\begin{itemize}[topsep=0pt]
    \item We present Group Relative Policy Optimization (GRPO), a novel and resource-efficient reinforcement learning method. Unlike Proximal Policy Optimization (PPO), GRPO eliminates the necessity for a critic network, instead leveraging group scores to estimate baselines, thereby substantially cutting training demands.
    \item Our experiments show that applying GRPO to our instruction-tuned model, DeepSeekMath-Instruct, yields notable performance improvements by utilizing only instruction-tuning datasets. In addition, we detect measurable gains in out-of-domain capabilities throughout reinforcement learning.
    \item We establish a coherent conceptual framework for situating RFT, DPO, PPO, and GRPO within a unified perspective. Through comprehensive empirical studies—including contrasts between online and offline learning, outcome versus process supervision, and comparisons of single-turn versus iterative reinforcement learning—we rigorously analyze core aspects of this unified approach.
    \item Guided by this integrated paradigm, we analyze the mechanisms underpinning the success of reinforcement learning and articulate several promising avenues for more effective application to LLMs.
\end{itemize}

\subsection{Summary of Evaluations and Metrics}

\begin{itemize}[topsep=0pt]
    \item \textbf{English and Chinese Mathematical Reasoning}: We undertake detailed evaluations on English and Chinese datasets spanning mathematical challenges from elementary to university levels. English assessments draw upon datasets such as GSM8K \citep{gsm8k}, MATH \citep{MATH}, SAT \citep{llemma}, OCW Courses \citep{minerva}, and MMLU-STEM \citep{mmlu}. Chinese assessments include MGSM-zh \citep{mgsm}, CMATH \citep{wei2023cmath}, Gaokao-MathCloze \citep{agieval}, and Gaokao-MathQA \citep{agieval}. Evaluations encompass the models' proficiency in composing independent text-based solutions as well as in solving problems programmatically via Python.

    On English tasks, DeepSeekMath-Base matches the closed-source Minerva 540B \citep{minerva} and outperforms all available open-source base models (including Mistral 7B \citep{mistral} and Llemma-34B \citep{llemma}), often by substantial margins, irrespective of pre-training regime. Remarkably, DeepSeekMath-Base achieves superior outcomes on Chinese mathematical benchmarks. We attribute this, in part, to our departure from prior practice \citep{minerva,llemma} of restricting math pre-training corpora to English, as our approach integrates high-quality non-English data as well. With further mathematical instruction tuning and reinforcement learning, our DeepSeekMath-Instruct and DeepSeekMath-RL models exhibit robust performance, marking the first occasion an open-source system attains over 50\% accuracy on the MATH competition dataset.
\end{itemize}

\item \textbf{Formal Mathematics}: DeepSeekMath-Base is assessed on the informal-to-formal theorem proving task introduced in \citep{dsp_proof}, utilizing the miniF2F benchmark \citep{minif2f} and employing Isabelle \citep{isabelle} as the proof assistant. The results indicate that DeepSeekMath-Base achieves robust performance in few-shot autoformalization scenarios.
\item \textbf{Natural Language Understanding, Reasoning, and Code}: To thoroughly evaluate general language comprehension, reasoning, and programming abilities, DeepSeekMath-Base is benchmarked on Massive Multitask Language Understanding (MMLU) \citep{mmlu}, which includes 57 diverse multiple-choice tasks, BIG-Bench Hard (BBH) \citep{bbh}, comprising 23 tasks that often demand multi-step reasoning, as well as the widely adopted code evaluation benchmarks HumanEval \citep{codex} and MBPP \citep{mbpp}. Pre-training on mathematical data enhances both linguistic understanding and reasoning proficiency.

\section{Math Pre-Training}

\subsection{Data Collection and Decontamination} This section details the methodology for assembling the DeepSeekMath Corpus using resources from Common Crawl. As illustrated in Figure \ref{fig:data_collect}, an iterative procedure is presented, systematically expanding a large-scale mathematical corpus by starting from a seed dataset—specifically, a limited but highly curated math-related collection. Importantly, this framework is versatile and can be extended to domains such as programming.

The initial step employs OpenWebMath \citep{openwebmath}, an authoritative set of high-quality mathematics web content, as the foundational seed corpus. Leveraging this dataset, a fastText model \citep{joulin2016fasttext} is trained to retrieve web pages similar to those in OpenWebMath. For training, 500,000 samples are drawn at random from the seed data as positive examples, alongside 500,000 negative examples selected from Common Crawl. Training is conducted using a publicly available library\footnote{\url{https://fasttext.cc}}, specifying a vector dimension of 256, a learning rate of 0.1, a maximum word n-gram length of 3, a minimum word occurrence threshold of 3, and 3 epochs. To further refine the input data, URL-based deduplication as well as near-duplicate removal are applied to the Common Crawl, reducing it to 40 billion HTML pages. Mathematical web content is subsequently extracted from this deduplicated set using the trained fastText model. To ensure high quality, the recalled pages are ordered by the model’s predicted scores, and only those with the highest rankings are retained. The amount of data preserved is further examined via pre-training on the top 40B, 80B, 120B, and 160B tokens. In the initial round, the top 40B tokens are selected for retention.

Following the initial round of data gathering, a substantial number of mathematical web pages remain uncollected. This limitation arises primarily from the fastText model's exposure to a set of positive examples that lacks breadth. To address this, we broaden the seed corpus by incorporating further mathematical web resources, thereby enhancing the diversity of the training data for the fastText model. Our procedure begins by partitioning the full Common Crawl dataset into non-overlapping domains, where a domain is composed of web pages that share the same base URL. For each domain, we determine the proportion of pages captured during the first round of collection. Domains in which more than 10\% of their web pages have been gathered are labeled as math-related (e.g., \url{mathoverflow.net}).
We then conduct a manual review to identify URLs linked to mathematical content within these selected domains (for instance, \url{mathoverflow.net/questions}).
Uncollected pages corresponding to these math-oriented URLs are subsequently integrated into the seed corpus. This process expands our set of positive samples and permits the training of a more robust fastText model, thereby enhancing the retrieval of mathematical data in subsequent rounds. After completing four iterations, our final dataset comprises 35.5M mathematical web pages, accumulating 120B tokens. Analysis after the fourth iteration indicates that approximately 98\% of the data had already been collected by the third round, prompting us to discontinue further collection.

To mitigate any risk of benchmark leakage, we apply the methodology from \cite{deepseek-coder} to exclude any web pages that include content from well-known English mathematical benchmarks (such as GSM8K \citep{gsm8k} and MATH \citep{MATH}) or Chinese benchmarks (including CMATH \citep{wei2023cmath} and AGIEval \citep{agieval}). Specifically, any text segment containing a 10-gram phrase that matches exactly with any substring present in the evaluation benchmarks is filtered out from our mathematical training dataset. In cases where the benchmark sample contains fewer than 10 but at least 3 grams, we enforce an exact match to remove any potentially contaminated web content.

\subsection{Evaluating the DeepSeekMath~Corpus Quality}

To assess the value of the DeepSeekMath~Corpus, we carry out pre-training studies comparing it against contemporary math-specific corpora, including:

\begin{itemize}[topsep=0pt]
    \item \textbf{MathPile} \citep{mathpile}: A composite corpus containing 8.9B tokens collected from sources such as textbooks, Wikipedia, ProofWiki, CommonCrawl, StackExchange, and arXiv, with the latter contributing more than 85\% of the total;
    \item \textbf{OpenWebMath} \citep{openwebmath}: A collection curated from CommonCrawl, filtered to emphasize mathematical content and totaling 13.6B tokens;
    \item \textbf{Proof-Pile-2} \citep{llemma}: An amalgamation comprising OpenWebMath, AlgebraicStack (10.3B tokens of mathematical programming data), and arXiv papers (28.0B tokens). In our experiments on Proof-Pile-2, we adhere to the arXiv:Web:Code mixing ratio of 2:4:1 as specified by \cite{llemma}.
\end{itemize}

\subsubsection{Training Setting}
\label{sec:quality-policy}
We fine-tune a general pre-trained language model consisting of 1.3B parameters—the same architecture as the DeepSeek LLMs \citep{deepseek-llm}, referred to as DeepSeek-LLM 1.3B—on various mathematical corpora. For each corpus, separate training is performed for a total of 150B tokens. All training runs utilize the HAI-LLM \citep{haillm} training infrastructure, selected for its efficiency and lightweight design. Consistent with DeepSeek LLMs' methodology, optimization employs the AdamW optimizer \citep{adamW} with settings $\beta_1=0.9$, $\beta_2=0.95$, and $\mathrm{weight\_decay}=0.1$. The learning rate follows a multi-stage scheduling strategy: it ramps up to the maximum after 2,000 warmup steps, drops to 31.6\% of its peak at the 80\% training mark, and declines further to 10.0\% of its peak after 90\% of training. The maximum learning rate is set to 5.3e-4. Training proceeds with a batch size of 4M tokens and a context window of 4K tokens.

\subsubsection{Evaluation Results}

\textbf{The DeepSeekMath~Corpus stands out for its superior quality, extensive multilingual coverage, and unparalleled scale.}

\begin{itemize}[topsep=0pt]
    \item \textbf{High-quality}:
    The downstream performance of models is measured on eight mathematical benchmark datasets utilizing few-shot chain-of-thought prompting \cite{cot}. According to Table \ref{tab:corpora_comparison}, models trained with DeepSeekMath~Corpus consistently outperform others. Figure \ref{fig:corpora_comparisons} illustrates that, after 50B training tokens (equivalent to one epoch over Proof-Pile-2), the DeepSeekMath-trained model still surpasses the Proof-Pile-2-trained counterpart, suggesting superior average data quality in DeepSeekMath.
    \item \textbf{Multilingual}:
    DeepSeekMath~Corpus incorporates mathematical content across multiple languages, with English and Chinese representing the majority. Table \ref{tab:corpora_comparison} demonstrates that training with DeepSeekMath~Corpus leads to notable enhancements in mathematical reasoning tasks in both English and Chinese. By comparison, prior mathematical corpora focus predominantly on English, providing marginal or even negative gains for Chinese mathematical problem-solving.
    \item \textbf{Large-scale}:
    In terms of size, DeepSeekMath~Corpus greatly exceeds existing mathematical datasets. As indicated by Figure \ref{fig:corpora_comparisons}, DeepSeek-LLM 1.3B attains steeper and more sustained progress in learning curves when trained on DeepSeekMath~Corpus. Meanwhile, alternative corpora, being much smaller, must be cycled through repeatedly during training, which leads to the model’s performance plateauing rapidly.
\end{itemize}

\subsection{Training and Evaluating DeepSeekMath-Base 7B}

This section presents DeepSeekMath-Base 7B, a foundational model designed with robust mathematical reasoning capabilities. The model's initialization leverages DeepSeek-Coder-Base-v1.5 7B \citep{deepseek-coder} and is subsequently trained over a corpus totaling 500B tokens. The data utilized during training is apportioned as follows: 56\% from the DeepSeekMath Corpus, 4\% sourced from AlgebraicStack, 10\% derived from arXiv, 20\% consisting of code from Github, and the remaining 10\% comprising natural language materials from the Common Crawl dataset in both English and Chinese. The training procedure generally follows the approach described in Section \ref{sec:quality-policy}, with notable modifications including setting the peak learning rate to 4.2e-4 and employing a batch size of 10M tokens.

To rigorously measure DeepSeekMath-Base 7B’s mathematical proficiency, we perform extensive evaluations. These include its capacity to independently generate complete mathematical solutions, its competency in utilizing auxiliary tools to tackle mathematical problems, and its ability to perform formal theorem proving. Furthermore, we explore the general aptitude of the base model beyond mathematics, covering its skills in natural language understanding, logical reasoning, and computer programming.

\paragraph{Mathematical Problem Solving with Step-by-Step Reasoning}

DeepSeekMath-Base’s aptitude for mathematical problem solving is assessed using few-shot chain-of-thought prompting \citep{cot} across eight distinct benchmarks in both English and Chinese. The selected benchmarks include tasks requiring quantitative reasoning (such as GSM8K \citep{gsm8k}, MATH \citep{MATH}, and CMATH \citep{wei2023cmath}) as well as multiple-choice items (including MMLU-STEM \citep{mmlu} and Gaokao-MathQA \citep{agieval}), representing a diverse array of mathematical domains spanning from elementary topics to university-level challenges.

As presented in Table \ref{tab:base_math_cot}, among open-source base models—including the commonly employed Mistral 7B \citep{mistral} and the newly introduced Llemma 34B \citep{llemma}, which was trained with Proof-Pile-2 \citep{llemma} for mathematical tasks—DeepSeekMath-Base 7B achieves the highest results across all eight evaluated benchmarks. On the challenging competition-grade MATH dataset, DeepSeekMath-Base attains a margin exceeding 10\% absolute improvement relative to all previous open-source base models, and even outperforms Minerva 540B \citep{minerva}, a proprietary model constructed atop PaLM \citep{palm} that is 77 times larger and fine-tuned with mathematical literature.

\paragraph{Mathematical Problem Solving with Tool Use} To assess program-assisted mathematical reasoning, we employ few-shot program-of-thought prompting \citep{pot,pal} on the GSM8K and MATH datasets. In this setup, models are guided to address each problem by composing a Python script, leveraging modules such as \textit{math} and \textit{sympy} for advanced calculations. The produced answer is then taken as the output of the executed code. Table \ref{tab:base_math_pot_proof} illustrates that DeepSeekMath-Base 7B establishes a new performance high, surpassing the prior best open-source model, Llemma 34B.

\paragraph{Formal Mathematics}

Automating formal proofs plays a crucial role in verifying the correctness and dependability of mathematical arguments and has attracted increasing research interest. We benchmark DeepSeekMath-Base 7B on the informal-to-formal proof generation task from \citep{dsp_proof}, where the objective is to create a formalized proof from an informal description, its corresponding formal representation, and an informal outline. Using miniF2F \citep{minif2f}, which evaluates formal Olympiad-level math, we prompt models in a few-shot manner to output formal proofs for Isabelle. Consistent with the protocol in \cite{dsp_proof}, the model outputs proof sketches, and the automated theorem prover Sledgehammer \citep{sledgehammer} completes the finer details. As reported in Table \ref{tab:base_math_pot_proof}, DeepSeekMath-Base 7B achieves strong results in automatic proof formalization.

\paragraph{Natural Language Understanding, Reasoning, and Code}

For assessing broader capabilities, we evaluate natural language understanding on the MMLU dataset \citep{mmlu}, reasoning proficiency on BBH \citep{bbh}, and coding skill using HumanEval \citep{codex} and MBPP \citep{mbpp}. Table \ref{tab:understanding_reasoning_code} reveals that DeepSeekMath-Base 7B produces marked improvements over DeepSeek-Coder-Base-v1.5 \citep{deepseek-coder} on both MMLU and BBH, demonstrating the effectiveness of mathematical training in advancing language comprehension and reasoning abilities. Moreover, by supplementing the training with code-related tokens, DeepSeekMath-Base 7B preserves the coding performance exhibited by DeepSeek-Coder-Base-v1.5 on HumanEval and MBPP. Overall, on these three reasoning and coding tasks, DeepSeekMath-Base 7B delivers substantial gains compared to the generalist Mistral 7B \citep{mistral}.

\section{Supervised Fine-Tuning}

\subsection{SFT Data Curation}

We assemble a comprehensive instruction-tuning dataset focused on mathematics, incorporating both English and Chinese problem statements spanning a wide spectrum of mathematical domains and difficulty. Each problem is matched with a corresponding solution rendered in various reasoning formats: chain-of-thought (CoT) \citep{cot}, program-of-thought (PoT) \citep{pot,pal}, and tool-integrated reasoning \citep{tora}. Altogether, the training corpus consists of 776,000 examples.

\begin{itemize}[topsep=0pt]
    \item \textbf{English mathematical datasets}: We provide tool-integrated annotations for GSM8K and MATH problems, and utilize a curated portion of MathInstruct \citep{MathInstruct} together with the training split from Lila-OOD \citep{lila}, which comprise problems solved via CoT or PoT approaches. The English dataset collection spans an array of mathematical subfields, including but not limited to algebra, probability, number theory, calculus, and geometry.
    \item \textbf{Chinese mathematical datasets}: For Chinese, we gather K-12 mathematics questions that cover 76 individual subtopics such as linear equations, each annotated with solutions in both chain-of-thought and tool-integrated reasoning formats.
\end{itemize}

\subsection{Training and Evaluating DeepSeekMath-Instruct 7B}

Here, we detail the supervised fine-tuning of DeepSeekMath-Instruct 7B, which builds upon the DeepSeekMath-Base foundation using our instruction-tuning corpus. During training, examples are concatenated at random until the sequence length reaches a cap of 4,000 tokens. The training process encompasses 500 iterations, utilizing a batch size of 256, and maintains a fixed learning rate of 5e-5 throughout.

We assess the mathematical proficiency of the models, both with and without tool augmentation, on four quantitative reasoning benchmarks available in English and Chinese. For comparative evaluation, we position our model alongside the strongest models of the current period:

\begin{itemize}[topsep=0pt]
    \item \textbf{Closed-source models} considered in our evaluation include: (1) the GPT series, with GPT-4 \citep{gpt4} and GPT-4 Code Interpreter~\footnote{\url{https://openai.com/blog/chatgpt-plugins\#\#code-interpreter}} representing the top-performing variants; (2) Gemini Ultra and Pro \citep{gemini}; (3) Inflection-2 \citep{inflection-2}; (4) Grok-1~\footnote{\url{https://x.ai/model-card}};
    along with the latest offerings from Chinese industry, including (5) Baichuan-3~\footnote{\url{https://www.baichuan-ai.com}}; and (6) the recently launched GLM-4~\footnote{\url{https://open.bigmodel.cn/dev/api\#glm-4}}
    
    from the GLM model series \citep{glm}.
\end{itemize}

These models are intended for broad application, the majority of which have been subjected to various alignment strategies.

\item \textbf{Open-source models} consist of: generic language models such as (1) DeepSeek-LLM-Chat 67B \citep{deepseek-llm}, (2) Qwen 72B \citep{qwen}, (3) SeaLLM-v2 7B \citep{seallm}, and (4) ChatGLM3 6B \citep{chatglm3}. Additionally, models specifically strengthened for mathematical reasoning are included, such as (5) InternLM2-Math 20B\footnote{\url{https://github.com/InternLM/InternLM-Math}}, which extends InternLM2 and undergoes dedicated mathematics pretraining and subsequent instruction tuning, (6) Math-Shepherd-Mistral 7B, utilizing PPO training \citep{schulman2017proximal} on Mistral 7B \citep{mistral} with a reward model guided by process supervision, (7) the WizardMath suite \citep{wizardmath}, which enhances the mathematical reasoning of both Mistral 7B and Llama-2 70B \citep{llama2} via evolve-instruct—an instruction-tuning method utilizing AI-generated instructions—paired with PPO training using mainly GSM8K and MATH datasets, (8) MetaMath 70B \citep{metamath}, which is Llama-2 70B fine-tuned on augmented GSM8K and MATH data, (9) ToRA 34B \cite{tora}, a version of CodeLlama 34B adapted for mathematical tool-based reasoning, and (10) MAmmoTH 70B \citep{MathInstruct}, where Llama-2 70B receives instruction tuning based on the MathInstruct collection.

Table \ref{tab:sft_rl_math} displays results for the evaluation scenario where the use of external tools is restricted. Here, DeepSeekMath-Instruct 7B excels at multi-step reasoning tasks, particularly outperforming all open-source counterparts and most proprietary systems (such as Inflection-2 and Gemini Pro) on the advanced MATH dataset by no less than 9\% in absolute terms. This holds true even when compared to considerably larger models (e.g., Qwen 72B) or those explicitly improved by math-oriented reinforcement learning methods (e.g., WizardMath-v1.1 7B). Although DeepSeekMath-Instruct achieves results on par with the proprietary Chinese models GLM-4 and Baichuan-3 on MATH, it has yet to reach the performance levels of GPT-4 and Gemini Ultra.

In settings where both natural language reasoning and code-based tools can be leveraged to address problems, DeepSeekMath-Instruct 7B attains nearly 60\% accuracy on MATH, outperforming all other open-source models reported. Across additional benchmarks, this model’s results are comparable to DeepSeek-LLM-Chat 67B, the previous state-of-the-art open-source model, despite being only one-tenth its size.

\section{Reinforcement Learning}

\subsection{Group Relative Policy Optimization}

Reinforcement learning (RL) is widely recognized as a means to further elevate the mathematical reasoning capacities of LLMs after Supervised Fine-Tuning (SFT) has been performed \citep{wang2023math,wizardmath}. This section details Group Relative Policy Optimization (GRPO), our proposed RL algorithm noted for its efficiency and effectiveness.

\subsubsection{From PPO to GRPO} 

Proximal Policy Optimization (PPO) \citep{schulman2017proximal}, an actor-critic reinforcement learning algorithm, is prevalently applied during the RL fine-tuning phase of LLM development \citep{ouyang2022training}. PPO updates the model by maximizing a clipped surrogate objective, as follows:

\begin{equation}
\footnotesize
    \mathcal{J}_{PPO}(\theta) = \mathbb{E}{[q \sim P(Q), o \sim \pi_{\theta_{old}}(O|q)]} \frac{1}{|o|} \sum_{t=1}^{|o|} \min \left[ \frac{\pi_\theta(o_{t} | q, o_{<t})}{\pi_{\theta_{old}}(o_{t} | q, o_{<t})} A_{t}, \text{clip} \left( \frac{\pi_\theta(o_{t} | q, o_{<t})}{\pi_{\theta_{old}}(o_{t} | q, o_{<t

In this context, $\pi_{\theta}$ and $\pi_{\theta_{old}}$ denote the present and prior policy models, respectively. The variables $q$ and $o$ represent questions and corresponding outputs, where the outputs are sampled from both the question dataset and the old policy $\pi_{\theta_{old}}$. The hyper-parameter $\varepsilon$, related to clipping, is introduced in PPO to enhance training stability. The term $A_t$ signifies the advantage, calculated using Generalized Advantage Estimation (GAE) \citep{gae}, relying on future rewards $\{r_{\ge t}\}$ and a trainable value function $V_{\psi}$. As a result, PPO necessitates concurrent training of a value function alongside the main policy model. To counteract potential overfitting of the reward model, it is common practice to incorporate a per-token KL divergence penalty between the policy and a reference model within the reward formulation, following the method outlined by \citet{ouyang2022training}:

\begin{equation}
    r_{t} = r_\varphi(q, o_{\le t}) - \beta \log\frac{\pi_{\theta}(o_{t}|q, o_{<t})}{\pi_{ref}(o_{t}|q, o_{<t})},
\label{eq:PPO-reward}
\end{equation}

In this equation, $r_\varphi$ stands for the reward model, $\pi_{ref}$ is a reference model—usually the initial SFT model—and $\beta$ is the scaling parameter for the KL penalty.

Since the value function used in PPO often matches the policy model in size, this setup imposes considerable demands on both memory and computation. Moreover, during RL optimization, the value function acts as a baseline to reduce variance when calculating the advantage. In large language models (LLMs), the reward model generally only provides a reward at the final token, complicating the training of a token-wise accurate value function. To address these issues, as illustrated in Figure \ref{fig:grpo}, we introduce Group Relative Policy Optimization (GRPO). GRPO eliminates the necessity for an extra value function as found in PPO. Instead, it uses the mean reward from a group of outputs, all generated in response to the same prompt, as a baseline. Specifically, for each question $q$, a batch of outputs $\{o_1, o_2, \cdots, o_G\}$ is generated by sampling from the previous policy $\pi_{\theta_{old}}$, and the optimization of the new policy aims to maximize:

\begin{equation}
\footnotesize
\begin{split}
    \mathcal{J}_{GRPO}(\theta) &= \mathbb{E}{[q \sim P(Q), \{o_i\}_{i=1}^G \sim \pi_{\theta_{old}}(O|q)]}  \\
    & \frac{1}{G}\sum_{i=1}^G\frac{1}{|o_i|} \sum_{t=1}^{|o_i|} \left\{ \min \left[ \frac{\pi_\theta(o_{i,t} | q, o_{i,<t})}{\pi_{\theta_{old}}(o_{i,t} | q, o_{i,<t})} \hat{A}_{i,t}, \text{clip} \left( \frac{\pi_\theta(o_{i,t} | q, o_{i,<t})}{\pi_{\theta_{old}}(o_{i,t} | q, o_{i,<t})}, 1 - \varepsilon, 1 + \varepsilon \right)  \hat{A}_{i,t} \right] - \beta \mathbb{D}_{KL}\left[\pi_{\theta} || \pi_{ref}\right]\right\} ,
\end{split}
\label{eq:GRPO-obj}
\end{equation}

Here, $\varepsilon$ and $\beta$ are the tuning parameters, and $\hat{A}_{i,t}$ represents the advantage determined by comparing the relative rewards among the members of each output group, a process to be expounded in the next subsections. The strategy of using group-wise relative advantage in GRPO is particularly well-suited to reward models, as such models are commonly trained using pairwise comparisons between outputs in response to the same question. Furthermore, rather than embedding the KL penalty directly into the reward as in PPO, GRPO applies regularization by augmenting the loss with the KL divergence between the new policy and the reference policy, thereby streamlining the computation of $\hat{A}_{i,t}$. In contrast to the KL penalty approach embedded in (\ref{eq

\begin{algorithm}[t]
  \small
  \caption{Iterative Group Relative Policy Optimization}
  \textbf{Input} initial policy model $\pi_{\theta_{\text{init}}}$; reward models $r_\varphi$; task prompts $\mathcal{D}$; 
  hyperparameters $\varepsilon$, $\beta$, $\mu$
  \begin{algorithmic}[1]
    \State Assign policy model $\pi_\theta \leftarrow \pi_{\theta_{\text{init}}}$
    \For{iteration = 1, \dots, I}
       \State Set reference model $\pi_{ref} \leftarrow \pi_{\theta}$
      \For{step = 1, \dots, M}
      \State Select a mini-batch $\mathcal{D}_b$ from $\mathcal{D}$
      \State Update previous policy model $\pi_{\theta_{old}} \leftarrow \pi_{\theta}$ 
      \State For each prompt $q \in \mathcal{D}_b$, generate $G$ outputs $\{o_i\}_{i=1}^G$ sampled from $\pi_{\theta_{old}} (\cdot \mid q) $
      \State Evaluate all outputs $o_i$ using $r_\varphi$ to obtain reward values $\{r_i\}_{i=1}^{G}$
      \State Estimate token-level group relative advantages $\hat{A}_{i,t}$ for each $t$-th token in $o_i$
      \For{GRPO iteration = 1, \dots, $\mu$}
        \State Optimize policy model $\pi_{\theta}$ with respect to the GRPO objective (Equation \ref{eq:GC-GRPO})
      \EndFor
    \EndFor \State Further train $r_\varphi$ using a replay-based continual learning method.
    \EndFor 
  \end{algorithmic}
  \textbf{Output} $\pi_\theta$
  \label{alg:iter-grpo}
\end{algorithm}

\subsubsection{Outcome Supervision RL with GRPO} Specifically, for any given question $q$, a collection of outputs $\{o_1, o_2, \cdots, o_G\}$ is generated by sampling from the policy before update, $\pi_{\theta_{old}}$. These outputs are then evaluated with a reward model to yield a corresponding set of $G$ reward scores, denoted $\mathbf{r}=\{r_1, r_2, \cdots, r_G\}$. The set of rewards is then standardized by subtracting the mean and dividing by the standard deviation of the group. In outcome supervision, the normalized score is attributed as the reward at the terminal token of each output $o_i$, and this same value is assigned as the advantage $\hat{A}_{i, t}$ for every token position in that output: $\hat{A}_{i, t} = \widetilde{r}_i = \frac{r_i- {\rm mean}(\mathbf{r})}{{\rm std}(\mathbf{r})}$. The policy is then updated via maximization of the target function specified in equation (\ref{eq:GRPO-obj}).

\subsubsection{Process Supervision RL with GRPO} Providing only a final reward per output, as in outcome supervision, may lack the granularity needed for effective feedback in intricate mathematical reasoning tasks. Building on the approach of \cite{wang2023math}, we consider process supervision, which offers feedback after each intermediate reasoning step. Formally, for each question $q$ and corresponding $G$ outputs $\{o_1, o_2, \cdots, o_G\}$, a process-oriented reward model scores each reasoning step. This produces nested reward sets: $\mathbf{R} = \{ \{r_1^{index(1)},\cdots,r_1^{index(K_1)}\}, \cdots,  \{r_G^{index(1)},\cdots,r_G^{index(K_G)}\} \}$, where $index(j)$ locates the terminal token of the $j$-th step, and $K_i$ records how many steps output $o_i$ contains. Each of these rewards is standardized by computing $\widetilde{r}_i^{index(j)} = \frac{r_i^{index(j)} - {\rm mean(\mathbf{R})}}{{\rm std(\mathbf{R})}}$. Next, the token-level advantage $\hat{A}_{i, t}$ is defined as the cumulative sum of all normalized future process rewards beginning from token $t$: $\hat{A}_{i, t} = \sum_{index(j)

\subsubsection{Iterative RL with GRPO} During reinforcement learning, reliance on an outdated reward model can become inadequate for guiding the evolving policy model. To address this, we implement an iterative approach for RL with GRPO. As detailed in Algorithm \ref{alg:iter-grpo}, each iteration involves constructing new datasets for the reward model by leveraging samples produced from the current policy model. We update the reward model iteratively, applying a replay strategy that includes 10\% previously collected data to maintain diversity. Subsequently, the updated policy model serves as the reference for further training, and we continue optimizing the policy model with feedback from the refreshed reward model.

\subsection{Training and Evaluating DeepSeekMath-RL}

Our reinforcement learning experiments utilize DeepSeekMath-Instruct 7B as the base policy model. RL training data comprises chain-of-thought-formatted problems drawn from GSM8K and MATH subsets within the SFT corpus, totaling about 144,000 questions. All other SFT data are withheld to focus on RL's influence in settings where benchmarks have minimal direct exposure during RL. Construction of the reward model's training corpus follows the procedure outlined in \citep{wang2023math}. The initial reward model is fine-tuned from DeepSeekMath-Base 7B using a 2e-5 learning rate. For GRPO-based training, the policy model's learning rate is set to 1e-6, and the KL penalty coefficient is maintained at 0.04. For every prompt, we sample $64$ completions, enforce a maximum sequence length of 1024 tokens, and set the batch size for training at 1024. After each exploration phase, only a single update step is performed on the policy model. The evaluation protocol for DeepSeekMath-RL 7B matches that of DeepSeekMath-Instruct 7B, categorizing GSM8K and MATH with chain-of-thought reasoning as in-domain, while all additional benchmarks are considered out-of-domain.

Table \ref{tab:sft_rl_math} presents the results of both open-source and proprietary models using chain-of-thought as well as tool-integrated reasoning approaches across English and Chinese evaluation sets. Our observations include: 1) DeepSeekMath-RL 7B achieves accuracies of 88.2\% on GSM8K and 51.7\% on MATH when employing chain-of-thought reasoning. This level of accuracy exceeds that of all open-source counterparts within the 7B–70B parameter range, and outperforms most closed-source alternatives as well. 2) Notably, DeepSeekMath-RL 7B's training involved only chain-of-thought instruction tuning data from GSM8K and MATH, beginning from DeepSeekMath-Instruct 7B. Although this training data is rather narrowly scoped, DeepSeekMath-RL 7B delivers consistently superior performance compared to DeepSeekMath-Instruct 7B across all reported benchmarks, emphasizing the advantages conferred by reinforcement learning.

\section{Discussion}
Here, we discuss insights derived from our pre-training and RL studies.

\subsection{Lessons Learnt in Pre-Training}

We first detail the lessons gained during pre-training. Unless stated otherwise, our methodology aligns with the training procedures described in Section \ref{sec:quality-policy}. For clarity, references to the DeepSeekMath Corpus pertain to an 89-billion-token dataset assembled in the second phase of data curation.

\subsubsection{Code Training Benefits Mathematical Reasoning}

There is a widely held, though largely untested, assumption that training on code enhances reasoning ability. Our findings provide evidence supporting this idea in the context of mathematics: code-centric training enhances a model's capacity for mathematical reasoning, regardless of the involvement of external tools.

To examine the influence of code training on mathematical reasoning, we evaluated both two-stage and single-stage training approaches, defined as follows:

\textbf{Two-Stage Training}

\begin{itemize}[topsep=0pt]
    \item \textbf{Code Training for 400B Tokens $\rightarrow$ Math Training for 150B Tokens}: DeepSeek-LLM 1.3B is pre-trained with 400B tokens of code, followed by an additional 150B tokens focused on mathematics.
    \item \textbf{General Training for 400B Tokens $\rightarrow$ Math Training for 150B Tokens}: For comparison, we conduct an experiment where the initial training stage uses 400B tokens sampled from a large general-purpose corpus (prepared by DeepSeek-AI) rather than code. This setting serves to evaluate whether exposure to code confers an advantage in enhancing mathematical reasoning capabilities compared to general data.
\end{itemize}

\textbf{One-Stage Training}

\begin{itemize}[topsep=0pt]
    \item \textbf{Math Training for 150B Tokens}: Here, DeepSeek-LLM 1.3B is directly trained using 150B tokens of mathematics data.
    \item \textbf{Training on a mixture of 400B Code Tokens and 150B Math Tokens}: We observe that transitioning to math training after code pre-training can reduce coding proficiency. Thus, we explore whether simultaneously training on a blend of code (400B tokens) and math (150B tokens) in a single stage can not only support mathematical reasoning but also reduce the likelihood of catastrophic forgetting.
\end{itemize}

\paragraph{Results}
Empirical results, detailed in Table \ref{tab:code-math} and Table \ref{tab:code-math-general-eval}, summarize downstream task performance for each training paradigm.

Engaging in code pre-training facilitates the development of program-assisted mathematical reasoning skills under both the sequential (two-stage) and integrated (one-stage) frameworks. As evidenced by Table \ref{tab:code-math}, dedicating the initial phase to code data in the two-stage regime substantially strengthens the model’s capability to resolve Python-based challenges such as GSM8K and MATH. Introducing subsequent math-specific training brings additional performance gains. Notably, within the one-stage approach, a combined regimen of code and math tokens effectively addresses the catastrophic forgetting seen in two-stage schedules, and concurrently fosters both coding competence (as shown in Table \ref{tab:code-math-general-eval}) and program-augmented mathematical problem solving (as indicated in Table \ref{tab:code-math}).

Furthermore, code-focused training alone is advantageous even for mathematical reasoning tasks not reliant on tools. In the two-stage setup, initial code training imparts a measurable boost, and subsequently, math-focused training capitalizes on this foundation to achieve peak performance. In contrast, blending code and math tokens within a single-stage pipeline impairs math reasoning without the aid of tools. We hypothesize that this outcome reflects the capacity constraints of DeepSeek-LLM 1.3B, which may limit its ability to concurrently acquire both coding and mathematical proficiencies.

\subsubsection{ArXiv Papers Appear Unhelpful for Enhancing Mathematical Reasoning}

ArXiv papers are routinely incorporated into pre-training datasets for mathematics-focused language models \citep{minerva,gpt-f,llemma,mathpile}. Nonetheless, thorough assessments of their influence on mathematical reasoning proficiency remain scarce. Somewhat unexpectedly, our empirical findings indicate that arXiv papers offer little advantage in fostering mathematical reasoning skills. To investigate this, we employ models spanning different parameter counts—including DeepSeek-LLM 1.3B and DeepSeek-Coder-Base-v1.5 7B \citep{deepseek-coder}—with arXiv datasets subjected to distinct preprocessing approaches:

\begin{itemize}[topsep=0pt]
    \item \textbf{MathPile} \citep{mathpile}: An 8.9B-token dataset constructed via heuristic-based cleaning and filtering, with over 85\% comprised of scientific arXiv articles;
    \item \textbf{ArXiv-RedPajama} \citep{redpajama}: The full collection of arXiv LaTeX files, cleaned of preambles, comments, macros, and bibliographies, totaling 28.0B tokens.
\end{itemize}

For evaluation, we train DeepSeek-LLM 1.3B on 150B tokens and DeepSeek-Coder-Base-v1.5 7B on 40B tokens for each arXiv dataset. Across the board, models trained exclusively on arXiv sources fail to demonstrate significant gains—and sometimes even exhibit performance declines—across a range of mathematical reasoning benchmarks. These include quantitative reasoning tasks like GSM8K and MATH (Table \ref{tab:arxiv-cot}), multiple-choice evaluations such as MMLU-STEM (Table \ref{tab:arxiv-cot}), and benchmarks targeting formal mathematics such as miniF2F (Table \ref{tab:arxiv_atp}).

That said, there are caveats to this result which restrict its generality. Specifically, we did not investigate:

\begin{itemize}[topsep=0pt]
    \item The contribution of arXiv data to math-related objectives outside the present scope, for example, the process of translating formal theorems or proofs into informal language;
    \item The interaction effects of arXiv data when combined with other sources during training;
    \item The possibility that arXiv content might prove more valuable at larger model scales than those examined here.
\end{itemize}

Accordingly, we identify these aspects as open directions for future research.

\subsection{Observations on Reinforcement Learning}
\subsubsection{In Pursuit of a Unified Framework} This section introduces a unified approach for comparing various training algorithms—including SFT, RFT, DPO, PPO, and GRPO—and presents experimental investigations into the determinants within this unified framework. Broadly, the gradient with respect to model parameter $\theta$ for a given training algorithm can be expressed as:

\begin{equation}
    \nabla_{\theta}\mathcal{J}_{\textcolor{red}{\mathcal{A}}}(\theta) = \mathbb{E}[\underbrace{(q,o) \sim \textcolor{red}{\mathcal{D}}}_{Data \ Source}]\left( \frac{1}{|o|} \sum_{t=1}^{|o|}  \underbrace{GC_{{\mathcal{A}}}(q, o, t, \textcolor{red}{\pi_{{rf}}})}_{Gradient \ Coefficient}  \nabla_{\theta}\log \pi_{\theta}(o_t | q, o_{<t})\right).
\label{eq:objective}
\end{equation}

This formulation is composed of three principal elements: (1) \textit{Data Source $\mathcal{D}$}, specifying the dataset utilized for training; (2) \textit{Reward Function $\pi_{{rf}}$}, providing the feedback signal that guides optimization; and (3) \textit{Algorithm $\mathcal{A}$}, responsible for transforming both the training samples and reward information into the gradient coefficient $GC$, thereby dictating the extent of reward or penalty during updates. We discuss several prototypical algorithms within this generalized framework:

\begin{itemize}
    \item  \textbf{Supervised Fine-tuning (SFT)}: The SFT method adapts a pretrained model through additional training on SFT data that have been selected by humans.
    \item \textbf{Rejection Sampling Fine-tuning (RFT)}: RFT continues training the SFT model by selecting a subset of its generated outputs—those verified to be correct—using SFT prompts, and employing these filtered samples for fine-tuning.
    \item \textbf{Direct Preference Optimization (DPO)}: DPO further optimizes the SFT model by training on pairs of augmented responses, which are produced by the SFT model and refined according to the DPO pairwise loss criterion.
    \item \textbf{Online Rejection Sampling Fine-tuning (Online RFT)}: Unlike standard RFT, Online RFT updates the policy model (initialized from the SFT model) by sampling and fine-tuning on outputs generated dynamically by the current policy model itself.
    \item \textbf{PPO/GRPO}: These methods start from an SFT model and update the policy by reinforcing it using samples drawn from its own latest version during training.
\end{itemize}

The elements comprising these methods are outlined in Table \ref{tab:data_policy}. For a comprehensive account of the derivation process, please consult Appendix \ref{app:analysis-rl}.

\paragraph{Observation about Data Source} We categorize the data sources into online sampling and offline sampling. Online sampling involves generating training data from the live exploration trajectories of the currently trained policy model. In contrast, offline sampling leverages training data collected from the outputs of the initial SFT model. Both RFT and DPO utilize an offline sampling approach, while Online RFT and GRPO adopt online sampling. 

According to Figure \ref{fig:rl-analysis}, Online RFT consistently surpasses RFT across two evaluation benchmarks. In the early phases of training, Online RFT's performance is on par with that of RFT; however, as training advances, Online RFT gains a clear and decisive lead, underscoring the effectiveness of online learning. This trend is anticipated: at the outset, the actor and the SFT model are highly similar, yielding only minor distinctions in the sampled data. Over time, the divergence in outputs between the actor and SFT becomes more pronounced, making real-time data acquisition increasingly beneficial.

\paragraph{Observation about Gradient Coefficient} The algorithm translates the reward signal into a gradient coefficient used to adjust model parameters. In our experiments, the reward function is instantiated either as a `Rule', which assesses the quality of a response based on answer accuracy, or as a `Model', in which a reward model is trained to assign scores to each response using rule-based judgments as training data. As demonstrated by Equations \ref{eq:GC-RFT} and \ref{eq:GC-GRPO}, there is a crucial distinction between GRPO and Online RFT: GRPO modifies its gradient coefficient dynamically in response to the reward value generated by the reward model, facilitating nuanced positive or negative reinforcement based on response quality. Conversely, Online RFT does not feature this adaptation; it fails to penalize incorrect outputs and instead reinforces all correctly answered responses equally, without regard to the relative merit indicated by the reward signal.

As illustrated in Figure \ref{fig:rl-analysis}, GRPO outperforms online RFT, underscoring the effectiveness of modifying the coefficients associated with positive and negative gradients. Moreover, GRPO+PS achieves better results than GRPO+OS, which suggests that utilizing more granular, step-sensitive gradient coefficients is advantageous. We also investigate iterative RL by running two cycles in our experiments. The findings in Figure \ref{fig:iter-rl} reveal that applying iterative RL notably boosts performance, especially during the first cycle.

\subsubsection{Why RL Works?} In this work, we apply reinforcement learning to a selected portion of the instruction tuning dataset and observe considerable performance gains relative to the instruction tuning baseline. To further elucidate the mechanisms behind the efficacy of reinforcement learning, we assess both Pass@K and Maj@K metrics for the Instruct and RL-enhanced models across two benchmarks. Figure \ref{fig:maj-pass} demonstrates that RL particularly improves Maj@K accuracy, while Pass@K remains unaffected. This pattern suggests that RL primarily strengthens the model's output distribution by making it more robust—put differently, \textbf{the gains appear to come from elevating the rank of the correct response among the TopK outputs rather than from increasing core competency.} Analogously, \citep{wang2023making} report a \textbf{misalignment problem} in SFT models on reasoning tasks and show that reasoning accuracy can be advanced through various preference alignment approaches \citep{yuan2023rrhf,song2023preference,wang2023making}.

\subsubsection{How to Achieve More Effective RL?}
\label{sec:effec_rl}
Our results confirm that RL yields substantial improvements on mathematical reasoning challenges. Additionally, we introduce a comprehensive framework to interpret several notable training approaches. Under this framework, all considered methods can be viewed as forms of RL, either in their full or reduced variants. As described in Equation \ref{eq:objective}, three central elements are identified: Data Source, Algorithm, and Reward Function. We offer several prospective research avenues related to each component.

\paragraph{Data Source} The data source is fundamental to any training strategy. For RL, we define the data source as consisting of unlabeled problems paired with output sequences generated by the policy model. In this study, we restrict the data to the instruction tuning questions and utilize straightforward nucleus sampling for response generation. We hypothesize that this limitation may account for the observation that RL in our pipeline predominantly improves Maj@K scores. Future work will investigate extending the RL framework to encompass out-of-distribution questions, and incorporate \textbf{advanced sampling (decoding) strategies}—for instance, those leveraging tree-search algorithms \citep{yao2023tree}. Furthermore, recent progress in \textbf{efficient inference techniques} \citep{xia-etal-2023-speculative,leviathan2023fast,kwon2023efficient,xia2024unlocking}, which significantly impact how effectively policy models explore, should also be considered essential.

\paragraph{Algorithms} Algorithms utilize both data and the reward signal to compute gradient coefficients for updating model parameters. According to Equation \ref{eq:objective}, current approaches effectively \textbf{TRUST} the reward function's feedback to modulate the conditional probability assigned to specific tokens, either increasing or decreasing it. Nevertheless, it remains infeasible to guarantee the reward signal’s fidelity at all times, particularly in highly intricate tasks. For instance, even the PRM800K dataset \citep{lightman2023let}, annotated by expert annotators with great care, contains about 20\% annotation errors\footnote{\url{https://github.com/openai/prm800k/issues/12\#issuecomment-1728491852}}. Consequently, we aim to investigate reinforcement learning algorithms that exhibit resilience to noisy or imperfect reward signals. Our hypothesis is that alignment techniques characterized as \textbf{WEAK-TO-STRONG} \citep{burns2023weak} could provoke foundational shifts in how learning algorithms are developed.

\paragraph{Reward Function} The reward function serves as the principal conduit for transmitting training signals. In reinforcement learning, this function is typically instantiated as a neural reward model. We identify three significant avenues for advancing reward models: 1) \textbf{How to improve the generalization capability of reward models.} The reward model should robustly generalize to novel, out-of-distribution queries and sophisticated decoding outcomes; otherwise, reinforcement learning may simply preserve the LLM's existing distribution instead of fundamentally enhancing its abilities; 2) \textbf{How to quantify and utilize the reward model's uncertainty.} Accurately capturing this uncertainty could create a valuable connection between limited reward models and weak-to-strong learning approaches; 3) \textbf{How to construct high-quality process reward models in an efficient manner} so as to offer detailed and nuanced training feedback throughout the reasoning pipeline \citep{lightman2023let,wang2023math}.

\section{Conclusion, Limitation, and Future Work}

This paper introduces DeepSeekMath, an open-source model that surpasses existing open-source systems on the MATH competition benchmark and nearly matches the results of proprietary counterparts. DeepSeekMath~begins with DeepSeek-Coder-v1.5 7B as its foundation and is continuously trained on a corpus of 500B tokens, including 120B mathematics-focused tokens gathered from Common Crawl. Through a detailed ablation study, we find that mathematical content extracted from the web holds considerable promise for high-quality data generation, while arXiv provides less utility than initially anticipated. We propose Group Relative Policy Optimization (GRPO), a modification of Proximal Policy Optimization (PPO), which substantially boosts mathematical reasoning performance while maintaining a lower memory footprint. Experimental results confirm that GRPO enhances capabilities even at the advanced performance level of DeepSeekMath-Instruct 7B. Furthermore, we propose a consolidated framework to contextualize various methods and outline several avenues for advancing reinforcement learning techniques.

Despite DeepSeekMath’s strong results on quantitative benchmarks, its performance in domains such as geometry and theorem-proving lags behind that of closed-source models. For example, preliminary testing reveals that the model is unable to solve triangle and ellipse-related questions, which may reflect a data selection imbalance during the pre-training and fine-tuning phases. Model scale also poses a limitation: DeepSeekMath’s performance in few-shot settings trails behind GPT-4, which exhibits marked improvements when provided with few-shot input, whereas DeepSeekMath’s few-shot and zero-shot outcomes are nearly identical. Looking forward, we aim to enhance our engineered data selection process to assemble richer pre-training datasets and pursue research outlined in Section \ref{sec:effec_rl} to further advance the reinforcement learning of LLMs.

\end{document}